\documentclass[11pt]{article}
\usepackage{amsmath,amssymb,amsthm,hyperref}
\newtheorem{theorem}{Theorem}
\newtheorem{proposition}{Proposition}
\title{Bayesian Neural Posterior Interfaces:\\Exact Closure and a Sharp Finite-Horizon Boundary}
\author{Contextual Composition research programme}
\date{13 November 2026}
\begin{document}\maketitle
\begin{abstract}
A Bayesian predictor is not just a next-output distribution: it is a stateful experiment whose evidence changes the next predictor. We give an exact criterion for a finite linear summary of a posterior to support this operation. The smallest stable summary is the likelihood-generated algebra, equivalently posterior masses of observationally identical parameter classes. A one-neuron model then supplies an arbitrarily long hierarchy: two priors can agree on every experiment of length at most $n-1$ yet disagree at length $n$. This is an elementary finite Bayesian inference example with written proofs and rational checks, not a new general stochastic contextad construction, practical neural-training method, or novelty claim.
\end{abstract}
\section{The operation that an interface must support}
Fix a nonempty finite parameter set $\Theta$, finite query set $X$, and finite output sets $Y_x$. For each $x,y$, a likelihood function
\[\ell_{x,y}:\Theta\longrightarrow[0,1],\qquad \sum_{y\in Y_x}\ell_{x,y}=1\]
is specified. A parameter is drawn once and retained throughout a run. Outputs of successive queries are conditionally independent given that parameter and the chosen queries. A query may depend on earlier observed outputs, but not on the hidden parameter. This includes a finite-prior Bayesian neural network: the weights are $\theta$ and the network specifies the output probabilities. The finite set is an exact modelling assumption, not an approximation theorem for continuous weights.

A posterior $q$ answers query $x$ with probability $z=q(\ell_{x,y})$, where $q(f)=\sum_\theta q(\theta)f(\theta)$. On observing $y$ with $z>0$, it becomes
\[ B_{x,y}(q)(f)=\frac{q(\ell_{x,y}f)}{q(\ell_{x,y})}. \tag{1}\]
No posterior is assigned to a zero-probability event. The normalized transition is partial; the unnormalized update $q\mapsto q\ell$ is total. This separation prevents impossible histories from acquiring arbitrary inferential content.

Let $S\subseteq\mathbb R^\Theta$ be a linear subspace containing $1$ and every likelihood. An interface retains $q|_S$, that is, the expectations of a basis of $S$. It therefore retains the probabilities of all next observations. The stronger requirement is that after any supported observation it can compute its own updated state, without returning to the discarded posterior.
\begin{theorem}[Exact posterior-summary closure]\label{thm:closure}
The following are equivalent.
\begin{enumerate}
\item For all priors $q,r$ with $q|_S=r|_S$, every observation of positive common probability gives $B_{x,y}(q)|_S=B_{x,y}(r)|_S$.
\item For every likelihood $\ell$ and $f\in S$, the pointwise product $\ell f$ belongs to $S$.
\end{enumerate}
When these conditions hold, the interface determines the law of every finite adaptive query/output transcript and can update exactly along it.
\end{theorem}
\begin{proof}
If products are in $S$, the numerator and denominator in (1) are determined by the interface. This proves the first implication and an explicit update algorithm. Induction on the transcript length proves equality of laws for any deterministic adaptive policy. A randomized policy using independent external randomness follows by averaging over that randomness.

Conversely, suppose $g=\ell f\notin S$. Finite-dimensional linear algebra gives a signed vector $v$ annihilating $S$ with $v(g)\ne0$. Since $1\in S$, $\sum v=0$. Start with the uniform positive prior $q_0$ and choose $\epsilon>0$ small enough that $q_\pm=q_0\pm\epsilon v$ are positive priors. They agree on $S$, hence on $\ell$. The likelihood is not identically zero (otherwise $g=0\in S$), so their common normalizer is positive. Their updated expectations of $f$ differ by $2\epsilon v(g)/q_0(\ell)$, contradicting the first condition.
\end{proof}

\section{The smallest stable interface and its exact meaning}
Let $A$ be the unital algebra generated by all likelihoods, with pointwise addition and multiplication. Define $\theta\sim\theta'$ when every likelihood has the same value at these parameters. These are observational parameter classes, not necessarily syntactic weight-permutation or other architectural symmetry classes.
\begin{theorem}[Minimal likelihood interface]\label{thm:minimal}
$A$ consists exactly of all functions constant on $\sim$-classes. It is the smallest linear interface containing $1$ and the likelihoods which is closed under Bayes updating. Its expectations are equivalent to the posterior mass of each class. Two priors give the same law for every finite adaptive experiment if and only if these class masses agree.
\end{theorem}
\begin{proof}
Products of likelihoods stay constant on classes, so $A$ is contained in the class-constant functions. For each class $C$ and each other class $D$, choose a likelihood $\ell_D$ separating representative points $c,d$. Then
\[e_C(\theta)=\prod_{D\ne C}\frac{\ell_D(\theta)-\ell_D(d)}{\ell_D(c)-\ell_D(d)}\]
is exactly the indicator of $C$ on the finite parameter set: it equals one on $C$ and has a zero factor on every other class. Thus all class indicators lie in $A$. Repeated multiplication shows any stable $S$ must contain $A$, while $A$ itself is stable.

Equal class masses give equal products of likelihoods and hence equal transcript probabilities. Conversely every monomial in likelihoods is the probability function of a prescribed sequence of queries and particular output word (repetition allowed). Equality of all these word probabilities implies equality of expectations of their linear span $A$, including the class indicators. Adaptivity is therefore not necessary for distinguishing different class masses, although the same interface supports it.
\end{proof}
This is minimality among the stated linear-expectation interfaces and an exact observational-equivalence result. It is not a lower bound for arbitrary encodings, approximate prediction, or Bayesian models restricted to one prescribed conjugate family. For $d$ distinct classes there are $d-1$ free posterior-mass coordinates. Hidden-unit permutations which preserve every allowed likelihood may legitimately be merged; equality only of the current mean prediction is much weaker.

One constructive algorithm starts with the span of $1$ and the likelihood vectors, repeatedly adds products of each basis vector with each likelihood, and stops when rank stabilizes. There are at most $|\Theta|$ dimensions. Rational likelihoods give exact rational linear algebra, though the parameter set can already be impractically large.

\section{A neural example with an arbitrarily late failure}
For $n\ge2$, let a single ReLU unit have weight $w_i=i/n$, $0\le i\le n$, input $x\in[0,1]$, and output probability
\[p_i(x)=\operatorname{ReLU}(w_i x)=w_i x.\]
The observed output is Bernoulli with this probability. Restrict allowed inputs to any finite subset of $[0,1]$ containing $1$. The hidden weight persists; new Bernoulli noise is drawn each call. This is a deliberately small Bayesian neural model, not a claim that a one-neuron architecture is itself difficult.

Define priors $q_E,q_O$ by normalized binomial mass on even and odd indices respectively:
\[q_E(i)=\begin{cases}\binom ni/2^{n-1}&i\text{ even},\\0&i\text{ odd},\end{cases}
\quad q_O(i)=\begin{cases}0&i\text{ even},\\\binom ni/2^{n-1}&i\text{ odd}.\end{cases}\]
\begin{proposition}[Sharp horizon separation]\label{prop:horizon}
These priors give identical laws for every adaptive experiment using at most $n-1$ calls. For $n$ calls at input $1$, their probabilities of all successes differ in absolute value by
\[\frac{n!}{2^{n-1}n^n}>0.\]
Consequently, after the common positive-probability history of $n-1$ successes at input $1$, their next-success predictions differ.
\end{proposition}
\begin{proof}
Each forward difference lowers the degree of a nonconstant polynomial by one and multiplies its leading coefficient by its previous degree. Expanding the $n$th iterated difference gives $\Delta^n P(0)=\sum_i(-1)^{n-i}\binom ni P(i)$. Thus it annihilates degrees below $n$ and gives $n!$ on $i^n$. Therefore
\[\sum_{i=0}^n(-1)^i\binom ni i^k=0\ (k<n),\qquad
\sum_{i=0}^n(-1)^i\binom ni i^n=(-1)^n n!.\]
The priors have equal moments through degree $n-1$ and the claimed difference at degree $n$. For a fixed transcript of an adaptive policy, its chosen inputs are fixed numbers, so its probability as a function of $w$ is a product of $k$ factors of the form $wx$ or $1-wx$. It has degree at most $k$. Equal moments imply equal transcript laws up to $n-1$ calls. Taking all successes at $x=1$ gives the displayed difference. Both priors have positive mass at some positive weight, so the shorter all-success history has common positive probability. Divide the two different length-$n$ probabilities by that common probability.
\end{proof}
For $n=2$, $q_E$ is half mass each at $w=0,1$, while $q_O$ is concentrated at $w=1/2$. Both initially predict success probability $1/2$ at input $1$. After one success, they predict $1$ and $1/2$ respectively. The significance is not just a correlation counterexample: no fixed finite horizon of tests certifies indefinite posterior composability uniformly over the family of finite models. The explicit separation shrinks with $n$; this is an exact, not quantitatively robust, lower bound.

For this model, retaining moments $m_j=q(w^j)$ for $0\le j\le k$ permits a query followed by $k-1$-moment output:
\[m'_j=\frac{x m_{j+1}}{x m_1}\quad(y=1),\qquad
m'_j=\frac{m_j-xm_{j+1}}{1-xm_1}\quad(y=0),\quad j<k.\]
Use only supported branches. Thus a horizon-indexed interface consumes one unit of remaining horizon at each observation. Trying to reuse the same truncated interface indefinitely is the error. On $n+1$ distinct weights, powers $1,w,\ldots,w^n$ span all functions by Vandermonde interpolation, and higher powers reduce modulo $\prod_{i=0}^n(t-i/n)$. The full finite moment interface closes exactly; discarded higher moments are not generally recoverable below this threshold.

\section{Placement in Contextual Composition}
The catalogue's Bayesian/stochastic neural direction asks for stochastic structure together with inference. Here a context is an available posterior interface (and, for truncations, a remaining horizon); content is a query followed by an output-dependent continuation. The evidence update changes the context used by that continuation. The square ``summarize then update'' versus ``update then summarize'' commutes exactly under Theorem~\ref{thm:closure}. Otherwise even though the next-output law factors through the summary, continuation semantics does not. The horizon-indexed formulas show a useful typed alternative to falsely claiming a persistent interface.

The retained realization can be written as the weighted joint experiment on $(\theta,y_1,\ldots,y_k)$, followed by projection to observable transcripts. Equal class-mass interfaces preserve all such finite projected laws. This is a finite operational example consistent with the programme's retained-witness versus law-quotient distinction; we do not derive new O1 interaction/coherence data or assert that every posterior-summary map is a native cell in an unspecified doctrine. Zeros are handled explicitly here; the inherited probability/List positive-support theorem is not being extended silently.

The existing consistent-random-function and checkpointing examples show lost correlations under forgetting or resampling. The extra contribution of this example is the necessary-and-sufficient posterior-update closure test, its constructive minimal interface, and the sharp finite-horizon hierarchy. The speculative-decoding example compares local masking with global conditioning; this one asks which stored inference state remains sufficient under repeated evidence. Neither comparison supports a novelty claim for the underlying probability theory.

\section{Checks, limits, and sources}
The accompanying standard-library Python script uses exact fractions. It checks the binomial hierarchy, every binary output word up to its boundary, genuinely adaptive inputs, moment-versus-posterior updating, rank growth, and a two-query model with a duplicate parameter class. These finite diagnostics supplement the written proofs. No Lean replay, numerical neural training, continuous-weight theorem, asymptotic approximation guarantee, or physical application validation is claimed.

Bayesian weight uncertainty and posterior predictive averaging are standard; see Blundell et al., \emph{Weight Uncertainty in Neural Network}, ICML 2015, \url{https://proceedings.mlr.press/v37/blundell15.html}. That work is not the source of the particular finite-interface theorem stated here. The algebra argument is elementary finite-dimensional linear algebra and polynomial interpolation; the horizon witness is the elementary finite-difference identity proved above. This note does not establish literature novelty.

Programme sources inspected: \emph{Integrated Relative Computation Acceptance Dossier} (finite probability control and retained semantics), the existing \emph{Constrained Speculative Sampling} example, and the catalogue's Bayesian/stochastic neural direction. A full new stochastic enrichment remains outside this note's scope.
\end{document}
